\chapter{Interest Groups, Bureaucracy and Regulation}\label{ch:b3:interest-groups-and-regulation}

Voters choose governments rarely and know little about most of what they do. Between
elections, policy is shaped by those who have much at stake in it and organise to
say so, and it is carried out by agencies whose staff have aims of their own. This
chapter builds the models of that process --- collective action, \omterm{def:b3:interest-groups-and-regulation:rent-seeking}{rent seeking},
regulation as a market for support, the bureau --- and tests them where the record
allows: occupational licences, trucking, \omterm{def:b3:interest-groups-and-regulation:privatisation}{privatisation} and \omterm{def:b3:interest-groups-and-regulation:privatisation}{contracting out}.

\begin{openingchart}{Share of U.S. workers who need a licence to do their job, by three
measures (percent): occupations licensed by State laws (early 1950s, drawn at the
upper bound of 5~\%, and 2008); a commissioned survey counting licences from every
level of government (2008); and the Current Population Survey's licence question,
asked since 2015. The measures differ; the break is drawn as a gap.}{Treasury, Council
of Economic Advisers and Department of Labor (2015), from Kleiner and Krueger; U.S.
Bureau of Labor Statistics, Current Population Survey, Table 49}
\begin{tikzpicture}
\begin{axis}[omtimesten, width=0.7\linewidth, height=0.4\linewidth, xmin=1945,
  xmax=2030, ymin=0, ymax=35, ylabel={Percent of workers}, legend pos=north west,
  legend style={font=\footnotesize}]
  \addplot[only marks, mark=*, omDef] table[x=year, y=share] {figdata/bachelor-3/22-interest-groups-and-regulation/out/lic_state.csv};
  \addplot[only marks, mark=square*, omThm] table[x=year, y=share] {figdata/bachelor-3/22-interest-groups-and-regulation/out/lic_survey.csv};
  \addplot[black, thick, mark=triangle*] table[x=year, y=share] {figdata/bachelor-3/22-interest-groups-and-regulation/out/lic_cps.csv};
  \legend{State licensing laws, Survey (all levels), Current Population Survey}
\end{axis}
\end{tikzpicture}
\end{openingchart}

% ledger: b3ig:wh
% ledger: b3ig:bls15
% ledger: b3ig:bls25
Fewer than one American worker in twenty needed a State licence in the early 1950s; by
2008 one in four did, and 29~\% counting local and Federal licences. The official survey,
which asks the question since 2015, puts the share between 21.6 and 22.4~\%, a different
measure on a flat path. Most of the growth came from new occupations being licensed
rather than from workers moving into old ones: changes in the mix of jobs explain
only a little more than a third of the rise from the 1960s to 2008. Each new licence was voted by a
legislature; this chapter asks who asked for them.

\begin{recall}
The grades 10--12 volume named \omterm{def:b3:interest-groups-and-regulation:rent-seeking}{rent seeking} and \omterm{def:b3:interest-groups-and-regulation:capture}{regulatory capture}; they are
re-founded here as models. The public-good contribution game is the Year 2
volume's, the deadweight loss of monopoly the Year 1 volume's, and the median voter
\cref{ch:b3:public-choice-voting}'s.
\end{recall}

\begin{notation}[Symbols of this chapter]\label{not:b3:interest-groups-and-regulation:symbols}
$V$ is a rent at stake; $x_i$ a player's spending to win it; $n$ the number of players or
members; $B(Q)$ and $C(Q)$ a bureau's total benefit and total cost.
\end{notation}

\section{The logic of collective action}\label{sec:b3:interest-groups-and-regulation:olson}

\begin{definition}[Collective action problem]\label{def:b3:interest-groups-and-regulation:collective-action}
\index{collective action problem}\index{concentrated benefits}\index{diffuse costs}
A \emph{collective action problem} is the under-provision of an effort that benefits
a whole group because each member bears the full cost of her own effort and gets
only a share of its benefit. A policy with \emph{concentrated benefits and diffuse
costs} gives a large gain to each member of a small group and a small loss to each
member of a large one.
\end{definition}

\begin{proposition}[Group effort falls with group size]\label{prop:b3:interest-groups-and-regulation:olson}
A group of $n$ identical members shares equally a benefit $v(G)$ of its total effort $G = \sum_i g_i$,
with $v$ increasing and strictly concave, and each member bears her own effort. In the
Nash equilibrium, $v'(G) = n$; so $G$ falls as $n$ rises, while the efficient effort solves
$v'(G) = 1$ whatever $n$.
\end{proposition}

\begin{proof}
Member $i$ chooses $g_i$ to maximise $v(G)/n - g_i$; the first-order condition is $v'(G)/n = 1$. Since
$v'$ is decreasing, the $G$ that solves $v'(G) = n$ is smaller the larger $n$. The group's total
payoff $v(G) - G$ is maximised at $v'(G) = 1$, which is above the equilibrium $G$ for every
$n > 1$.
\end{proof}

A few producers each with much to gain outspend millions of consumers each with
little to lose, even when the consumers' total loss is larger. Licensing fits: the
members of a trade are few and know each other, and their customers are many and
scattered.

\section{Rent seeking and its cost}\label{sec:b3:interest-groups-and-regulation:rent-seeking}

\begin{definition}[Rent seeking; rent dissipation]\label{def:b3:interest-groups-and-regulation:rent-seeking}
\index{rent seeking}\index{rent dissipation}\index{Tullock contest}
\emph{Rent seeking} is the spending of resources to obtain a transfer created by
policy --- a licence, a tariff, a monopoly --- rather than to produce. \emph{Rent
dissipation} is the share of the rent consumed by that spending; in a \emph{Tullock
contest} each player wins the rent with a probability equal to her share of total
spending.
\end{definition}

\begin{proposition}[Dissipation in a symmetric contest]\label{prop:b3:interest-groups-and-regulation:tullock}
With $n$ risk-neutral players competing for a rent $V$, each spends $(n - 1)V/n^2$, and total
spending is $(n - 1)V/n$.
\end{proposition}

\begin{proof}
Player $i$'s expected payoff is $Vx_i/X - x_i$ with $X = \sum_j x_j$. The first-order condition is
$V(X - x_i)/X^2 = 1$. In a symmetric equilibrium $x_i = x$ and $X = nx$, so $V(n - 1)x/(n^2x^2) = 1$, or $x = (n - 1)V/n^2$; the
payoff is concave in $x_i$, so this is a maximum. Total spending is $nx = (n - 1)V/n$.
\end{proof}

\begin{omchart}{Share of the rent dissipated in a symmetric Tullock contest, by number
of players (percent).}{Analytic: \cref{prop:b3:interest-groups-and-regulation:tullock}}
\begin{tikzpicture}
\begin{axis}[omrecord, width=0.56\linewidth, height=0.34\linewidth, xmin=0, xmax=21,
  ymin=0, ymax=100, xlabel={Number of players $n$}, ylabel={Percent of the rent}]
  \addplot[omDef, thick, mark=*, mark size=1.2pt] table[x=n, y=share] {figdata/bachelor-3/22-interest-groups-and-regulation/out/tullock.csv};
\end{axis}
\end{tikzpicture}
\end{omchart}

Two contestants burn half the rent; ten burn 90~\%. The cost of a monopoly created by
policy is then not only the deadweight triangle but much of the rectangle of
profit as well, spent on lobbying, lawyers and campaigns that produce nothing.

\begin{omfigure}
\begin{tikzpicture}[econ, x=0.55cm, y=0.55cm]
  \draw[econ axes] (0,0) -- (11,0) node[right, font=\scriptsize] {$Q$};
  \draw[econ axes] (0,0) -- (0,8) node[above, font=\scriptsize] {price};
  \draw[curve, omDef] (0,7.5) -- (9.5,0.85) node[above right, font=\scriptsize] {demand};
  \draw[curve] (0,2.5) -- (10.5,2.5) node[right, font=\scriptsize] {marginal cost};
  \fill[omDef!15] (0,2.5) rectangle (3.57,5);
  \fill[omThm!35] (3.57,5) -- (7.14,2.5) -- (3.57,2.5) -- cycle;
  \node[font=\scriptsize, align=center] at (1.8,3.75) {Tullock\\rectangle};
  \node[font=\scriptsize, anchor=west] at (7.3,3.6) {Harberger triangle};
  \draw[guide] (0,5) -- (3.57,5); \node[left, font=\scriptsize] at (0,5) {$p^m$};
  \draw[guide] (3.57,0) -- (3.57,2.5); \draw[guide] (7.14,0) -- (7.14,2.5);
  \node[below, font=\scriptsize] at (3.57,0) {$Q^m$};
  \node[below, font=\scriptsize] at (7.14,0) {$Q^c$};
\end{tikzpicture}
\par\smallskip{\small A monopoly created by policy. The triangle is the deadweight
loss of the restricted output; the rectangle, the monopoly profit, is a transfer if
nobody competes for it and a social cost to the extent that contestants spend it
away.}
\end{omfigure}

\section{Regulation as a market for support}\label{sec:b3:interest-groups-and-regulation:capture}

\begin{definition}[Regulatory capture; Peltzman model]\label{def:b3:interest-groups-and-regulation:capture}
\index{regulatory capture}\index{Peltzman model}
\emph{Regulatory capture} is the use of regulation for the benefit of the regulated
industry rather than its customers. The \emph{Peltzman model} is the model in which a
regulator sets a price to maximise political support $M(\pi, S)$, increasing in the
industry's profit $\pi$ and in consumers' surplus $S$.
\end{definition}

\begin{proposition}[The regulated price lies between]\label{prop:b3:interest-groups-and-regulation:peltzman}
Let profit $\pi(p)$ rise from zero at the competitive price $p^c$ to its maximum at the monopoly
price $p^m$, let consumers' surplus $S(p)$ fall in $p$, and let $M$ be differentiable with
$M_\pi > -M_S S'(p^c)/\pi'(p^c)$ at $p^c$ (support from the first dalen of profit outweighs the consumers'
loss). Then the support-maximising price lies strictly between $p^c$ and $p^m$.
\end{proposition}

\begin{proof}
Support along prices is $m(p) = M(\pi(p), S(p))$, with $m'(p) = M_\pi\pi'(p) + M_SS'(p)$. At $p^c$, the
assumption gives $m'(p^c) > 0$, so a slightly higher price raises support. At $p^m$, $\pi'(p^m) = 0$ and
$S'(p^m) < 0$, so $m'(p^m) = M_SS'(p^m) < 0$: a slightly lower price raises support. A maximum of
$m$ on $[p^c, p^m]$ therefore lies inside the interval.
\end{proof}

\begin{model}{The capture model of regulation}
\assumes Regulators trade the favours of regulation for political support, and
small organised groups supply support more cheaply than large ones.
\predicts Regulation raises the prices and incomes of concentrated producers above
competitive levels, and less so where consumers are organised.
\tested Licensed workers earn about 14~\% more than similar unlicensed ones, while
governmental \omterm{def:b3:interest-groups-and-regulation:licensing}{certification}, which does not bar the uncertified, adds much less; % ledger: b3ig:kk
across U.S. states, stricter dental licensing raised prices without improving the
dental health of recruits. % ledger: b3ig:kku
When trucking was deregulated, the union wage premium fell from 50~\% to under 30~\%:
union workers had captured more than two-thirds of the industry's regulatory
rents. % ledger: b3ig:rose
\breaks Some regulation serves large, diffuse groups against concentrated ones ---
rules on pollution are the standard case --- and some licences may protect buyers
who cannot judge quality before they buy.
\end{model}

\begin{definition}[Occupational licensing]\label{def:b3:interest-groups-and-regulation:licensing}
\index{occupational licensing}\index{certification}
\emph{Occupational licensing} is a law that allows only holders of a government
licence to be paid for work in an occupation; \emph{certification} lets anyone work
but reserves a title to those who pass a test.
\end{definition}

\begin{measure}[The welfare cost of licensing]
Comparing occupations that some U.S. states license and others do not, an
equilibrium model in which licensing restricts supply but may raise demand through
quality finds that licensing raises wages and hours and lowers employment. The
average welfare loss is 12~\% of the surplus of the occupation; workers bear 70~\% of it
and consumers 30~\%. Buyers' higher willingness to pay offsets 80~\% of the higher
prices, and higher wages repay workers for 60~\% of the training the licence
requires. % ledger: b3ig:ks
The estimate is silent on doctors and lawyers, who are licensed everywhere.
\end{measure}

\section{The budget-maximising bureau}\label{sec:b3:interest-groups-and-regulation:bureau}

\begin{definition}[Budget-maximising bureau (Niskanen model)]\label{def:b3:interest-groups-and-regulation:bureau}
\index{budget-maximising bureau}\index{Niskanen model}
A \emph{budget-maximising bureau} is an agency that supplies a service to a sponsor
--- a legislature --- in exchange for a budget, knows the cost of the service better
than the sponsor does, and chooses the output that maximises its budget subject to
the sponsor's accepting it (the \emph{Niskanen model}).
\end{definition}

\begin{proposition}[The bureau doubles output]\label{prop:b3:interest-groups-and-regulation:niskanen}
Let the sponsor's marginal benefit be $a - bQ$ and the marginal cost a constant $c$, with $c < a \le 2c$. The
efficient output is $Q^* = (a - c)/b$; a bureau that offers output $Q$ for a budget equal to the
total benefit $B(Q)$, subject to $B(Q) \ge C(Q)$, chooses $Q = 2Q^*$.
\end{proposition}

\begin{proof}
$B(Q) = aQ - \tfrac b2Q^2$ is the largest budget the sponsor accepts for $Q$; it rises with $Q$ up to
$a/b$. The bureau wants the largest $B(Q)$ it can cover, so it pushes $Q$ up until the
budget no longer covers cost: $aQ - \tfrac b2Q^2 = cQ$, or $Q = 2(a - c)/b = 2Q^*$, provided $2Q^* \le a/b$, that is
$a \le 2c$, so that marginal benefit is still positive there (otherwise it stops at $a/b$). At $2Q^*$, all the surplus is spent: the extra output
from $Q^*$ to $2Q^*$ costs more than it is worth by exactly the surplus earned up to $Q^*$.
\end{proof}

\begin{omfigure}
\begin{tikzpicture}[econ, x=0.55cm, y=0.55cm]
  \draw[econ axes] (0,0) -- (11,0) node[right, font=\scriptsize] {$Q$};
  \draw[econ axes] (0,0) -- (0,8) node[above, font=\scriptsize] {dalens per unit};
  \draw[curve, omDef] (0,7) -- (10,2) node[right, font=\scriptsize] {marginal benefit};
  \draw[curve] (0,4.5) -- (10.5,4.5) node[right, font=\scriptsize] {marginal cost};
  \fill[omDef!15] (0,4.5) -- (0,7) -- (5,4.5) -- cycle;
  \fill[omThm!35] (5,4.5) -- (10,4.5) -- (10,2) -- cycle;
  \draw[guide] (5,0) -- (5,4.5); \draw[guide] (10,0) -- (10,2);
  \node[below, font=\scriptsize] at (5,0) {$Q^*$};
  \node[below, font=\scriptsize] at (10,0) {$2Q^*$};
  \node[font=\scriptsize] at (1.6,5.3) {surplus};
  \node[font=\scriptsize] at (8.2,3.9) {loss};
\end{tikzpicture}
\par\smallskip{\small The Niskanen bureau. Output beyond $Q^*$ costs more than it is
worth; at $2Q^*$ the loss equals the surplus, and the sponsor is just willing to pay.}
\end{omfigure}

\begin{model}{The budget-maximising bureau}
\assumes The bureau alone knows its costs, the sponsor knows its own demand, and
the bureau's staff gain from a larger budget.
\predicts Public agencies produce too much, or at too high a cost; competition
from other suppliers lowers the cost of the same service.
\tested Weakly. International studies of \omterm{def:b3:interest-groups-and-regulation:privatisation}{contracting out} --- putting public services
out to competitive tender --- report savings in the order of 20~\% with no loss of
quality; % ledger: b3ig:dj
the savings fit the model, with a caveat: the services put out to tender may be
those whose quality is easiest to specify.
\breaks Sponsors monitor, compare agencies and cut budgets; staff may care about
the mission more than the budget; and contracting fails for services whose quality
is hard to write into a contract.
\end{model}

\section{Privatisation and deregulation}\label{sec:b3:interest-groups-and-regulation:privatisation}

\begin{definition}[Privatisation]\label{def:b3:interest-groups-and-regulation:privatisation}
\index{privatisation}\index{contracting out}
\emph{Privatisation} is the sale of a state-owned enterprise, or of a controlling
stake in it, to private owners; \emph{contracting out} keeps the service public and
buys it from private suppliers.
\end{definition}

\begin{record}[Privatisation in the record]
A survey of the empirical studies of \omterm{def:b3:interest-groups-and-regulation:privatisation}{privatisation}, in countries in transition from
central planning and elsewhere, finds that in most settings \omterm{def:b3:interest-groups-and-regulation:privatisation}{privatisation} works:
the firms become more efficient, more profitable and financially healthier, and
investors in the sales do well. % ledger: b3ig:mn
\emph{Grade 3}: a survey of studies that mostly compare firms before and after sale,
with the firms chosen for sale not at random.
Ownership is not the whole story, and the theory says why: a private monopoly sold without competition
or regulation keeps a monopoly's incentive to restrict output.
The best-measured comparison of the transition shows how much the setting matters.
Following every initially state-owned manufacturing firm in four countries, Brown,
Earle and Telegdy estimated that \omterm{def:b3:interest-groups-and-regulation:privatisation}{privatisation} raised productivity by 15~\% in
Romania, 8~\% in Hungary and 2~\% in Ukraine, and lowered it by 3~\% in Russia, where
gains from domestic owners appeared only five years after the sale; sales to
foreign owners raised it by 18--35~\% everywhere. % ledger: b3ig:bet.priv
\emph{Grade 2}: firm panels with controls for which firms were sold, to which the
estimates are sensitive.
\end{record}

\begin{definition}[Regulatory impact assessment]\label{def:b3:interest-groups-and-regulation:ria}
\index{regulatory impact assessment}
A \emph{regulatory impact assessment} is the cost-benefit analysis of a proposed
regulation, carried out by the government before it adopts the rule, with the
methods of \cref{ch:b3:cost-benefit-analysis}.
\end{definition}

The public-choice models suggest why such assessments exist and why they are hard:
the costs of a rule on an organised industry are counted by those who bear them,
the benefits to diffuse beneficiaries by nobody in particular. A required
assessment puts both on the same page before the vote.

\begin{record}[Regulatory cost estimates against outcomes]
Harrington, Morgenstern and Nelson compared the costs that \omterm{def:b3:interest-groups-and-regulation:ria}{regulatory impact
assessments} had forecast for 28 United States rules with the costs measured after
the rules took effect. The total cost had been overestimated for 14 rules and
underestimated for 3. The cost per unit of pollution abated was often right; the
errors came mostly from the quantity of abatement, which suggests that benefits were
overstated too, and rules that used prices or tradable permits were the ones whose
unit costs were overestimated most consistently. % ledger: b3ig:hmn.ria
Unlike the project forecasts of \cref{ch:b3:cost-benefit-analysis}, these errors lean towards overstating costs, through
baselines that were hard to measure, incomplete compliance and innovations nobody
foresaw.
\end{record}

\section{Exercises}\label{sec:b3:interest-groups-and-regulation:exercises}

\begin{exercise}[$\star$]\label{exo:b3:interest-groups-and-regulation:1}
From the opening chart, by how many points did the share of workers covered by
State licences rise from the early 1950s to 2008? Why can the 2015 survey value not be
compared with the 2008 survey value?
\end{exercise}

\begin{exercise}[$\star$]\label{exo:b3:interest-groups-and-regulation:2}
Four players compete in a symmetric \omterm{def:b3:interest-groups-and-regulation:rent-seeking}{Tullock contest} for a rent of 10 million dalens. How
much does each spend, and how much of the rent is dissipated?
\end{exercise}

\begin{exercise}[$\star$]\label{exo:b3:interest-groups-and-regulation:3}
Demand is $p = 20 - Q$ and marginal cost 4. A licence creates a monopoly. Find the Harberger
triangle and the Tullock rectangle.
\end{exercise}

\begin{exercise}[$\star$]\label{exo:b3:interest-groups-and-regulation:4}
What evidence would contradict the capture model in a licensed profession?
\end{exercise}

\begin{exercise}[$\star\star$]\label{exo:b3:interest-groups-and-regulation:5}
A sponsor's marginal benefit from a bureau's output is $100 - 2Q$ and marginal cost is 60. Find
the efficient output, the bureau's output and the deadweight loss.
\end{exercise}

\begin{exercise}[$\star\star$]\label{exo:b3:interest-groups-and-regulation:6}
A tariff gives 50 million dalens to 10 firms and costs 60 million to 2 million households.
With the collective-action proposition, explain which group lobbies harder. What
would change if the firms were 10{,}000 farms?
\end{exercise}

\begin{exercise}[$\star\star$]\label{exo:b3:interest-groups-and-regulation:7}
An invented regression of log hourly wages on a licence indicator and controls
gives a coefficient of 0.13 (s.e.\ 0.02). Translate it into a percentage, and say what
omitted variable could bias it.
\end{exercise}

\begin{exercise}[$\star\star$]\label{exo:b3:interest-groups-and-regulation:8}
In the \omterm{def:b3:interest-groups-and-regulation:capture}{Peltzman model}, demand for the industry's product rises. Using the
proposition, explain why the regulated price may rise while the regulator still
leaves the industry short of the monopoly profit.
\end{exercise}

\begin{exercise}[$\star\star$]\label{exo:b3:interest-groups-and-regulation:9}
What quality data would support licensing as consumer protection?
\end{exercise}

\begin{exercise}[$\star\star\star$]\label{exo:b3:interest-groups-and-regulation:10}
Two players with stakes $V_1$ and $V_2$ compete in a \omterm{def:b3:interest-groups-and-regulation:rent-seeking}{Tullock contest}. Find the equilibrium
spending of each and the total, and show that the total is less than either stake.
\end{exercise}

\begin{exercise}[$\star\star\star$]\label{exo:b3:interest-groups-and-regulation:11}
A natural monopoly with fixed cost 16, marginal cost 2 and demand $p = 12 - Q$ is
privatised under a price cap. Find the price that just covers cost and compare it
with the unregulated monopoly price. Why do regulators set caps from forecasts of
cost rather than from the firm's realised costs?
\end{exercise}

\begin{exercise}[$\star\star\star$]\label{exo:b3:interest-groups-and-regulation:12}
In invented data on 40 states, 2012--2023, fifteen states loosen a trade's licence in 2018.
Estimate the effect on the log price of the service by difference in differences
with state and year fixed effects, and compare it with the before--after change in
the reforming states.
\end{exercise}

\section{Problem: Brenland Licenses Its Electricians}\label{sec:b3:interest-groups-and-regulation:problem}

\begin{problem}[{Weekend problem --- Brenland licenses its electricians}]\label{pb:b3:interest-groups-and-regulation:electricians}
In Brenland, inverse demand for electricians' hours is $P = 100 - 0.05Q$ and supply is
$P = 20 + 0.03Q$, with $Q$ in thousands of hours a year and $P$ in dalens an hour. A proposed
licence would cap hours at 800 thousand a year. % figdata: electricians_stats

\medskip\noindent\textbf{Part I --- the market without a licence.}
\begin{enumerate}
\item Find the competitive quantity and price.
\item Compute consumers' surplus.
\item Compute producers' surplus.
\item Why is the market a candidate for licensing in the first place?
\item What would a licence have to achieve to raise total surplus?
\end{enumerate}

\medskip\noindent\textbf{Part II --- the licence.}
\begin{enumerate}[resume]
\item With the cap, find the price.
\item Find the supply price of the 800th thousand hours.
\item Compute the rent per hour and in total.
\item By how much does the price exceed the competitive price, in percent?
\item Compute the Harberger triangle.
\item Compute the consumers' loss.
\item Compute the producers' gain, and check that the consumers' loss equals it plus
the triangle.
\end{enumerate}

\medskip\noindent\textbf{Part III --- the contest.} The electricians' trade association
lobbies for the licence and a consumer group against it, in a \omterm{def:b3:interest-groups-and-regulation:rent-seeking}{Tullock contest}
fought each year, in which each side's stake is its gain or loss that year.
\begin{enumerate}[resume]
\item State each side's stake.
\item Find each side's spending.
\item Find the probability that the licence passes.
\item Compute total rent-seeking spending.
\item Which side spends more, and why does the side with the smaller stake win less
often?
\item Why would an unorganised consumer side spend far less than this?
\item What happens to the probability of passage if consumers do not organise at all?
\end{enumerate}

\medskip\noindent\textbf{Part IV --- the repeal.} Years later, four of Brenland's eight
regions repeal the licence. Mean prices (dalens an hour) were 60 before and 53 after in
the repealing regions, and 58 and 57 in the others.
\begin{enumerate}[resume]
\item Compute the before--after change in the repealing regions.
\item Compute the difference-in-differences estimate.
\item What assumption makes it the effect of the repeal?
\item Compare it with the model's prediction for a full repeal.
\item Give one reason the estimate might be smaller than predicted.
\item State the result: the total social cost of the licence, Harberger plus Tullock,
in millions of dalens.
\end{enumerate}
\end{problem}
